%!TEX root = ../dissertation.tex

\chapter{Conclusion}
\label{chp:conclusion}

This thesis aimed to address the problem of explainability applied to predictive process monitoring on object-centric event logs. As predictive process monitoring methods become more and more widespread, the reliance of these methods on black-box predictive models such as neural networks may be a detriment for user adoption since research has shown that process stakeholders will only adopt methodologies they feel they can trust. An explainability layer on top of these models can help create trust by providing insight into the elements driving a model's predictions. The proposal is a four-component explainability framework meant to provide stakeholders with clear and understandable insights into the predictive model being used. By providing both factual and counterfactual explanations, the stakeholder is able to understand which specific elements in their inputs are driving the model's prediction as well as how changing different elements in the input can lead to other desired outputs.


In order to test the framework two different synthetic OCEL 2.0 datasets are used, meant to emulate the processes of an Order Management and a Logistics company. For each dataset one heterogeneous graph neural network was trained per KPI, and two KPIs were chosen per dataset. The results show that, across the entire prefix set, the heterogeneous architecture is roughly on par with a homogeneous graph-based predictor, with overlapping confidence intervals across all four KPIs and no consistent winner. However, when compared only on last-event cases (which contain the greatest possible amount of information regarding a process execution) the picture becomes clearer. In these cases the heterogeneous model significantly outperformed every baseline on two of the four KPIs (Logistics' Depart and LoadToVehicle), remained the strongest choice but was statistically indistinguishable from the k-dim GNN baseline on a third (Order Management's PayOrder), and underperformed on the fourth (Order Management's PackageDelivered). This shows that the heterogeneous architecture is able to take advantage of richer input information as a process extends, though the benefit is KPI-dependent rather than universal.


The proposed explainability framework was then applied to all four KPI-specific regression models. The experimental evaluation showed that LOO analysis could be used as an efficient way of identifying impactful elements within the input graph, helping to curtail the number of elements for which the computationally costly Shapley Values analysis could be performed on. Similarly, gradient-based methods can be used to calculate the attribution value of element features in the graph and identify the most impactful elements on the model's prediction. Using these two identifier methods and applying Shapley Values to their chosen elements helped to provide a comprehensive set of input elements and the magnitude of their impact on the model's prediction. These results are then presented to the stakeholder in easy to understand graphs and tables that provide a clear look at the important elements, their rankings and their impact on the given prediction. In order to further create trust in the stakeholder, assessment metrics are also provided for the given factual explanations. Characterization and sparsity are shown to help the stakeholder understand how well the given explanation is at identifying the important elements for the model and separating them from the less important features.


Aside from giving insights into the particular elements driving the prediction, the framework is also able to provide the stakeholder with a counterfactual explanation method to help identify changes in the process that can lead to a desired change such as reducing the time remaining until an event occurs. By presenting these elements in a visually concise way, the stakeholder is able to identify key features in the process that can lead to reductions in the time it takes for a process to finish, or elements that are hindering the process as a whole.


Current limitations for the explainability framework pertain to the limited amount of explainability options available for use with heterogeneous models. While methods like HGExplainer and HTGExplainer exist, they could not be applied in this thesis given certain limitations (specifically for link-prediction/recommender-systems tasks and a failed port from DGL into PyTorch respectively). Additionally, while Shapley Values offer a comprehensive analysis of the features in a graph input and their impact on the prediction it also has a high computational cost which limits its possible applications as process executions grow in size and complexity. A further limitation observed during evaluation concerns the sensitivity of small temporal train/validation/test splits to whichever particular slice of a process's timeline each split happens to cover. For example, for one KPI (Order Management's "Orders to PayOrder"), the validation split's target values were found to be intrinsically less variable than the training split's, which on its own is enough to produce a pattern that could otherwise be mistaken for overfitting. Future work using a validation scheme less sensitive to which particular window of time it draws from, such as rolling-origin or k-fold temporal validation rather than a single fixed split, could mitigate this.


Future work in this space could seek to apply a similar explanation to a global, cross-trace view of the process rather than the current local-only implementation, following the work of de Leoni \& Volpato \cite{deleonieVolpato}. Similarly, a method seeking to use HTGExplainer could be constructed following a DGL-pipeline in order to validate its use on predictive process monitoring \cite{li2023}.


Ultimately, the proposed framework is able to provide insights into the predictive process monitoring model being applied to object-centric event logs. By providing stakeholders with explanations quantifying the impact of the most important elements of the inputs, metrics that assess the effectiveness of those explanations, and a counterfactual method to compare two distinct process executions, the framework provides a comprehensive look at the driving forces behind the predictive model.